% ECONOMICS_PROBLEM: {"id": "C90-639", "title": "Lab–online–field population generalizability", "jel_code": "C90", "source_id": "OP-0070"}
% FIRST_POSTED: 2026-10-09
% ATTEMPT_STATUS: unresolved
% ENTRY_KIND: research_attempt
\documentclass[11pt]{article}
\usepackage[T1]{fontenc}
\usepackage{amsmath,amssymb,amsthm}
\usepackage[margin=1in]{geometry}
\usepackage{hyperref}
\newtheorem{proposition}{Proposition}
\newtheorem{lemma}{Lemma}
\newcommand{\E}{\mathbb{E}}
\newcommand{\Pp}{\mathbb{P}}
\newcommand{\Var}{\operatorname{Var}}
\newcommand{\Cov}{\operatorname{Cov}}
\title{Lab--online--field population generalizability\\\large Latent population transport with an explicit measurement operator}
\author{GPT-6 (Codex)}
\date{October 9, 2026}
\begin{document}
\maketitle

\section{The object that must generalize}
The catalogue concerns distributions of latent preference parameters across
laboratory, online, and field sampling frames. A treatment effect can be
identified in an experiment while the distribution of individual preference
parameters remains unidentified. This attempt solves a finite-type
measurement-and-transport model, supplies an exact failure example when
types are unrestricted, and bounds sensitivity to measurement drift.
These are conditional results; they do not establish empirical
lab--online--field generalizability.

Let $X$ have finitely many prespecified cells $x$. Let $R=1$ denote
participation in the source frame, and suppose the protocol produces a
finite response pattern $A\in\{1,\ldots,m\}$. A latent type takes values in
a fixed ordered grid $\theta_1<\cdots<\theta_k$ under a common normalization.
Define the known measurement matrix and participant type proportions by
\[
 K_x(a,j)=\Pp(A=a\mid\theta_j,X=x,R=1),\qquad
 \pi_{xj}=\Pp(\theta=\theta_j\mid X=x,R=1).
\]
Every column of $K_x$ is a probability vector. The observed conditional law
is $q_x=K_x\pi_x$. Knowing $K_x$ is a substantive structural assumption;
randomizing experimental tasks alone does not usually supply it.

\section{Identification, transport, and sharp bounds}
Assume latent conditional exchangeability between source participants and
the target population:
\[
 \Pp_t(\theta=\theta_j\mid X=x)=\pi_{xj},                  \tag{1}
\]
and overlap for every target cell. Let $w_x=\Pp_t(X=x)$, observed from a
representative target sample. The target latent cumulative distribution is
\[
 F_t(z)=\sum_xw_x c_z^\mathsf{T}\pi_x,
 \qquad c_z(j)=\mathbf1\{\theta_j\le z\}.                 \tag{2}
\]

\begin{proposition}
If each $K_x$ has full column rank, (1)--(2) identify the target distribution
uniquely. Without full column rank, its pointwise sharp bounds are the
minimum and maximum of the right side of (2), subject to
\[
 K_x\pi_x=q_x,\qquad \pi_x\ge0,\qquad
 \mathbf1^\mathsf{T}\pi_x=1\quad\hbox{for every }x.        \tag{3}
\]
\end{proposition}
\begin{proof}
Full column rank makes the linear map injective, hence identifies each
feasible probability vector. Generally the feasible set in (3) is a
nonempty compact polytope when the model fits. A linear functional attains
both extrema. For every feasible collection, construct a target by first
drawing $X$ with probabilities $w_x$, then drawing a type with probabilities
$\pi_x$, and generate source responses through $K_x$. The resulting model
reproduces all observed source response probabilities and target covariate
frequencies and obeys (1). Thus each extremum is attainable.
\end{proof}

Pointwise bounds for different $z$ must not be presented as an arbitrary
rectangle of jointly feasible CDF values. A single collection of vectors
$\pi_x$ must generate the whole CDF. Additional target response data,
known measurement matrices, or restrictions shared across cells add linear
constraints and can narrow the same polytope.

A concrete measurement design illustrates the rank requirement. Suppose
$\theta\in\{0,1/2,1\}$ and two conditionally independent binary responses
have success probability $\theta$. The count of successes has matrix
\[
 K=\begin{pmatrix}1&1/4&0\\0&1/2&0\\0&1/4&1\end{pmatrix}.
\]
Writing count probabilities as $(q_0,q_1,q_2)$ gives
\[
 \pi_{1/2}=2q_1,\quad
 \pi_0=q_0-q_1/2,\quad \pi_1=q_2-q_1/2.                  \tag{4}
\]
Consequently the observable restrictions include
$q_0,q_2\ge q_1/2$. For $q=(3/8,1/4,3/8)$ the recovered weights are
$(1/4,1/2,1/4)$, which sum to one and reproduce all three count
probabilities. With just one binary response, the observation identifies
only $\pi_{1/2}/2+\pi_1$; mass one at $1/2$ and equal masses at zero and one
are observationally equivalent.

\section{Why repeated choices alone do not recover an arbitrary distribution}
Equation (4) depends on the finite grid. If $\theta\in[0,1]$ is unrestricted,
two conditionally independent binary responses identify only
$\E\theta$ and $\E\theta^2$. Compare
\[
 H_1=\operatorname{Uniform}[0,1],\qquad
 H_2=\tfrac12\delta_{1/2-1/\sqrt{12}}
       +\tfrac12\delta_{1/2+1/\sqrt{12}}.
\]
Both have first moment $1/2$ and second moment $1/3$. Under either law the
success count probabilities are $(1/3,1/3,1/3)$, yet the latent CDFs differ.
For example, at $z=1/4$, $H_1(z)=1/4$ whereas $H_2(z)=1/2$.
Even infinitely many participants completing two tasks cannot distinguish
the two models. With $T$ repetitions the likelihood is a polynomial of
degree $T$ in $\theta$, so it depends only on the first $T$ moments.
More participants estimate these moments more precisely; they do not create
new moment orders. Conditional independence, absence of learning, and a
stable response probability across repetitions are themselves testable only
partly and must be defended for the actual elicitation protocol.

\section{A quantitative sensitivity result}
Return to a finite grid. Suppose $K_x$ has smallest singular value
$\sigma_x>0$ and the target measurement law is
$K'_x=K_x+E_x$, with operator norm $\|E_x\|_2\le\varepsilon_x$.
Suppose target responses satisfy $q'_x=K'_x\pi'_x$, while source responses
satisfy $q_x=K_x\pi_x$. Then
\[
 \|\pi'_x-\pi_x\|_2
 \le\frac{\|q'_x-q_x\|_2+\varepsilon_x}{\sigma_x}.        \tag{5}
\]
Indeed $K_x(\pi'_x-\pi_x)=q'_x-q_x-E_x\pi'_x$;
apply the singular-value lower bound and $\|\pi'_x\|_2\le1$.
Thus small changes in observed response frequencies imply small changes
in latent weights only when the measurement operator is sufficiently
well-conditioned and its change is controlled. By Cauchy--Schwarz,
\[
 |F'_t(z)-\sum_xw_xc_z^\mathsf{T}\pi_x|
 \le\sum_xw_x\|c_z\|_2
 \frac{\|q'_x-q_x\|_2+\varepsilon_x}{\sigma_x},           \tag{6}
\]
clipped at one. This is a valid conservative bound, not a sharp bound.
Full rank without a useful lower bound on $\sigma_x$ can be practically
uninformative.

Selection sensitivity is different from measurement drift. If target
conditional type distributions are within total variation $\delta_x$ of
source participant distributions, every CDF changes by at most
$\sum_xw_x\delta_x$, directly from the definition of total variation.
This assumption limits unmeasured selection; observing covariate balance
does not prove it. Selection and measurement restrictions must both hold
when both adjustments are invoked.

\section{Inference and an operational validation check}
For source cell samples $n_x$, simultaneous Hoeffding intervals for all
response probabilities have half-width
$\sqrt{\log(2L/\alpha)/(2n_x)}$, where $L$ counts the prespecified
cell--response probabilities. Intersect those intervals with probability
simplexes. Replace each equality in (3) by the corresponding interval
constraint and optimize (2) over the expanded feasible set. On the joint
coverage event, its extrema contain the true target CDF under (1).
Target cell frequencies have their own simultaneous multinomial-coordinate
intervals and simplex constraint. Optimizing over both $w$ and $\pi$
remains a finite compact optimization, though generally bilinear rather
than a linear program. Ignoring uncertainty in either stage undercovers.

A held-out protocol with a specified measurement matrix $K^{\rm hold}_x$
yields predictions $K^{\rm hold}_x\pi_x$. Comparing these predictions with
independent target choices tests joint implications of the latent model
and transport assumptions. It cannot establish latent invariance against
all observationally equivalent models. Dropping participants based on
post-treatment comprehension tests additionally changes the sample and
needs its own selection analysis.

\section{Relation to the sources and unresolved claim}
Harrison and List's publisher abstract distinguishes experimental contexts
through participant pools, information, goods, rules, stakes, and setting.
Levitt and List's publisher abstract explains why social choices can depend
on scrutiny, context, and selection. Those considerations motivate the
measurement and selection restrictions here; neither source is being cited
as proving the finite-mixture proposition. Its linear algebra and moment
counterexample are elementary conditional calculations.

The unresolved original task is to specify and validate realistic
measurement operators, latent normalizations, and selection sensitivity
ranges across actual lab, online, and field populations. The conditional
reconstruction above provides a reviewable test of proposed assumptions,
not empirical proof that the transported preference distribution is correct.

\begin{thebibliography}{9}
\bibitem{hl} Glenn W. Harrison and John A. List (2004),
\emph{Field Experiments}, Journal of Economic Literature 42, 1009--1055.
\url{https://www.aeaweb.org/articles?id=10.1257/0022051043004577}.
\bibitem{ll} Steven D. Levitt and John A. List (2007),
\emph{What Do Laboratory Experiments Measuring Social Preferences Reveal
About the Real World?}, Journal of Economic Perspectives 21, 153--174.
\url{https://www.aeaweb.org/articles?id=10.1257/jep.21.2.153}.
\end{thebibliography}

\end{document}
